\section{引言}

实现可容错服务器的一种常见方法是primary/backup方法[1]，%
也就是始终准备一台backup服务器，在primary服务器故障(fail)时接管(take over)。%
backup服务器的状态必须始终与primary服务器几乎完全一致，%
这样当primary服务器failure时，backup服务器才能立即take over，%
并且以一种对外部客户端隐藏failure且不丢失数据的方式take over。%
复制backup服务器状态的一种方法，是几乎连续地把primary服务%
器所有状态的变化，包括CPU、内存和I/O设备的变化，%
都发送到backup服务器。然而，发送这些状态所需的带宽，%
尤其是内存变化所需的带宽，可能非常大。\emph{(译者注：这种方式被%
称为状态同步，常见的做法是先全量同步，后续同步“增量状态”，也就是上文
提到的“状态变化”)}%

另一种复制服务器的方法有时被称为状态机(state-machine)方法[13]，%
它可以使用少得多的带宽。其思想是把服务器建模为%
deterministic state-machine：让它们从相同的初始状态开始，%
并确保它们以相同顺序接收相同的输入请求，从而保持同步\emph{(译者注：%
这个思路就是复制状态机：给相同的初始态输入相同的操作，得出的%
结果一样，有点机械决定论那味了)}。%
由于大多数服务器或服务都有一些non-deterministic operation，%
\emph{(译者注：假设你有两台机器，都是单核心，同样的硬件配置，核心跑的同%
样的程序，然而它们也可能有不同的“输出”。我举个例子，你的代码里有“时间相关”%
的操作，比如读取CPU时钟，或者应用程序会读取gettimeofday，这样%
即使两个机器配置一样且同时开始运行，也会让“时间敏感”的程序输出不同；%
此外，在操作系统中，随机数也是这样，它实际上综合了来自操作系统的%
各类“混沌的信息”，这类混沌信息也许一部分来自“时间”，因此也会有不确定性)}%
因此必须使用额外的协调机制，确保primary服务器和backup服务器保持同步%
。不过，使primary服务器和backup服务器保持同步所需的额外信息量，%
远小于primary服务器中正在变化的状态量(主要是内存更新)。%

实现协调机制来确保物理服务器的deterministic execution%
[14]是困难的，尤其是在处理器频率不断提高时。相比之下，%
运行在hypervisor之上的VM是实现state-machine方法的极好平台。%
VM可以被视为一个定义良好的state-machine，%
其操作就是被虚拟化机器(包括所有设备)的操作。与物理服务器一样，%
VM也有一些non-deterministic operation(例如读取日时钟%
，或interrupt投递\emph{(译者注：对于两台机器，interrupt投递%
必须完全由primary控制时序，如果任由两台机器执行自己的中断，那么%
会由于中断插入的位置不同导致状态分岔。常见的请求就是两个机器，设置了10ms的时间中断，%
后续运行相同的代码，最后可能因为物理因素，比如芯片频率导致中断插入的位置不一样，%
进而使状态不一致，执行流分岔)}，因此必须向backup VM发送额外信息，%
确保它保持同步。由于hypervisor完全控制VM的执行，包括所有输入的投递，%
因此hypervisor能够捕获primary VM上%
non-deterministic operation的全部必要信息，%
并在backup VM上正确replay这些操作。%
\\\emph{译者注：hypervisor技术可以理解为将硬件拆分给多个虚拟机%
使用，虽然主板上只有一套物理硬件，但是可以依靠“多路复用”的思路实现拆分。%
我们可以选择直接在机器上部署裸金属hypervisor(不依赖于操作系统)。%
读者可以把这种裸金属hypervisor理解为一个用于管理VM虚拟化的简易操作系统，%
在裸机安装这个“微操作系统”后，可以配置它的网络(IP、Mask、Gateway，DNS)，%
外界可以直接进入它的Web UI进行管理，并且它提供了RESTful的API可以创建/销毁/%
暂停/快照虚拟机等功能，直接供外界调用，管理该机器的虚拟机资源。%
\\实际上本文所述的FT系统，需要的主要不是“虚拟化将硬件多路复用”的能力，
而是利用Hypervisor对VM执行边界的控制能力：截获输入、非确定性事件、虚拟中断、%
外部输出，并让backup按相同轨迹重放。实现这些能力需要修改%
hypervisor代码。}%

因此，state-machine方法可以在通用硬件上的VM中实现，不需要修改硬件，%
从而可以立即为最新的微处理器实现fault tolerance。此外，%
state-machine方法所需的低带宽，%
使primary和backup主机在物理上相隔更远成为可能。例如，%
backup VM可以运行在分布于一个园区内的物理机器上，这比运行在同一栋楼%
有更高的可靠性。%

我们在VMware vSphere 4.0平台上使用primary/backup方法%
实现了fault-tolerant VM。该平台以高效方式运行完全虚拟化的%
x86 VM。由于VMware vSphere实现了完整的%
x86 VM，我们能够自动为任何x86操作系统和应用提%
供fault tolerance能力。使我们能够记录primary VM的执%
行并确保backup VM完全相同执行的基础技术称为deterministic replay%
[15]。VMware vSphere Fault Tolerance(%
FT)基于deterministic replay，但增加了构建完整fault tolerance%
系统所需的额外协议和功能\emph{(译者注：后续会介绍为了对外表现为%
“未出现故障切换”，看起来具有输出一致性而设计的\textbf{Output Rule}；%
以及VM FT检测故障的方法；还有为了防止脑裂采取的额外措施；还有后面复杂的%
网络/磁盘IO故障切换后的一致性问题如何解决。所有的这些都需要进行讨论)}。%
除了提供硬件fault tolerance之外，%
我们的系统还会在failure后自动恢复冗余：%
它会在本地cluster中任意可用服务器上启动一个新的backup VM。目前，%
deterministic replay和VMware FT的生产版本都只支持单处理%
器VM。记录-重放多核VM的执行过程的版本仍处于开发阶段，%
它存在严重的性能问题，因为几乎每次对共享内存的访问都可能是不确定的。%

Bressoud和Schneider[3]描述了一个面向HP PA-RISC平台的%
fault-tolerant VM原型实现。我们的方法与之类似，%
但出于性能原因做出了一些根本改变，并考察了多种设计替代方案。此外，%
为了构建一个高效且可供运行企业应用的客户使用的完整系统，%
我们还必须设计并实现系统中的许多额外组件，并处理大量实际问题。%
与讨论过的大多数其他实际系统类似，我们只尝试处理fail-stop failure%
[12]，也就是那些能够在failed服务器造成错误的外部可见行为之前被检测到的%
服务器failure。%
\\\emph{译者注：fail-stop可以是机器简单地直接宕机，也可以是机器尚未%
立刻掉电、内部状态不明，但已经不再对外产生可观察输出。%
对于很多服务器程序来说，客户端只能通过网络观察服务器行为；%
在这一前提下，如果系统能够保证某一时刻之后该服务器不再向外%
发送网络报文、不再响应客户端请求，那么从客户端视角看，它就%
可以被抽象为一种合理的fail-stop故障。(不过，后续读者会%
看到，如果该服务器还能继续写共享存储、访问外部设备或产生其他
外部副作用，仅仅“不发网络包”还不够安全)}

本文其余部分组织如下。首先，我们描述基本设计，并详细说明我们的基础协议，%
这些协议确保在backup VM于primary VM failure后take over%
时不会丢失数据。然后，我们会详细描述为构建一个健壮、完整、%
自动化的系统，我们所必须解决的许多实际问题。我们还描述实现%
fault-tolerant VM时出现的若干设计选择，并讨论这些选择中的权衡。%

接下来，我们给出我们的实现针对若干benchmark和真实企业应用的性能结果。%

最后，我们描述相关工作并作总结。
